\section{Boundary Traversal of a Binary Tree}
\label{sec:trees-boundary}

\problemstatement{Given the root of a binary tree, return an
anti-clockwise boundary traversal: the left boundary (top-down, excluding
leaves), then all leaves left-to-right, then the right boundary
(bottom-up, excluding leaves). The root is included once.}

\begin{patternbox}
\emph{``Boundary,'' ``outline,'' ``perimeter''} decomposes into three
independent walks that must be stitched without double-counting: the left
edge going down, the leaves going across, and the right edge coming back
up. The only real care is avoiding overlap where these parts meet -- which
is why leaves are deliberately excluded from the two edges.
\end{patternbox}

\subsection*{Understanding the problem carefully}

The perimeter of the tree consists of three disjoint pieces. The
\textbf{left boundary} is the path you get by always going left (or right
if there is no left child), from the root's left child down, stopping
before leaves. The \textbf{leaves} are every node with no children, in
left-to-right order. The \textbf{right boundary} is the mirror -- always
go right (or left otherwise) -- collected then reversed so it reads
bottom-up. Excluding leaves from the two boundaries prevents a corner leaf
from being listed twice.

\begin{ideabox}[Three Walks, No Overlap]
(1) Add the root (unless it is itself a leaf). (2) Walk the left boundary
top-down, adding non-leaf nodes. (3) Add all leaves via any DFS,
left-to-right. (4) Walk the right boundary top-down into a temporary list,
adding non-leaf nodes, then append that list \emph{reversed}. Leaves are
skipped in steps~2 and~4 precisely so step~3 owns them exclusively.
\end{ideabox}

\subsection*{Algorithm}

\begin{enumerate}
  \item If root is not a leaf, add root's value.
  \item Left boundary: start at root.left; while the node exists, if it is
        not a leaf add it; move to its left child, or right child if no
        left.
  \item Leaves: DFS the whole tree; whenever a node has no children, add
        it.
  \item Right boundary: start at root.right; collect non-leaf nodes
        top-down (preferring right children) into a list; append that
        list reversed.
\end{enumerate}

\subsection*{Dry run}

\dryrun{Tree: root $1$; left $2$ (children $4$,$5$); right $3$ (right
child $6$, whose left child $7$).}

\begin{itemize}
  \item Root $1$ (not a leaf): add $1$.
  \item Left boundary from $2$: $2$ not a leaf, add $2$; go to $4$ -- a
        leaf, stop. Boundary so far: $1,2$.
  \item Leaves L$\to$R: $4,5,7$. Now: $1,2,4,5,7$.
  \item Right boundary from $3$: $3$ not leaf add $3$; go to $6$, not leaf
        add $6$; $7$ is a leaf, stop. List $[3,6]$ reversed $=[6,3]$.
  \item Result: $1,2,4,5,7,6,3$.
\end{itemize}

\subsection*{C++ implementation}

\begin{lstlisting}
#include <bits/stdc++.h>
using namespace std;

struct TreeNode {
    int val;
    TreeNode* left;
    TreeNode* right;
    TreeNode(int v) : val(v), left(nullptr), right(nullptr) {}
};

bool isLeaf(TreeNode* n) {
    return n->left == nullptr && n->right == nullptr;
}

void addLeaves(TreeNode* n, vector<int>& out) {
    if (n == nullptr) return;
    if (isLeaf(n)) { out.push_back(n->val); return; }
    addLeaves(n->left, out);
    addLeaves(n->right, out);
}

vector<int> boundaryTraversal(TreeNode* root) {
    vector<int> out;
    if (root == nullptr) return out;
    if (!isLeaf(root)) out.push_back(root->val);

    // left boundary (top-down, non-leaf)
    TreeNode* cur = root->left;
    while (cur) {
        if (!isLeaf(cur)) out.push_back(cur->val);
        cur = cur->left ? cur->left : cur->right;
    }

    // all leaves, left to right
    addLeaves(root, out);

    // right boundary (top-down into temp, then reversed)
    vector<int> tmp;
    cur = root->right;
    while (cur) {
        if (!isLeaf(cur)) tmp.push_back(cur->val);
        cur = cur->right ? cur->right : cur->left;
    }
    for (int i = (int)tmp.size() - 1; i >= 0; i--) out.push_back(tmp[i]);
    return out;
}

int main() {
    TreeNode* root = new TreeNode(1);
    root->left = new TreeNode(2);
    root->left->left = new TreeNode(4);
    root->left->right = new TreeNode(5);
    root->right = new TreeNode(3);
    root->right->right = new TreeNode(6);
    root->right->right->left = new TreeNode(7);

    for (int x : boundaryTraversal(root)) cout << x << " ";
    cout << "\n";
    return 0;
}
\end{lstlisting}

\begin{complexitybox}
\textbf{Time Complexity.} $O(n)$ -- the leaf DFS is $O(n)$ and each
boundary walk is $O(h)$. \textbf{Space Complexity.} $O(n)$ for the output
and $O(h)$ recursion for the leaf pass.
\end{complexitybox}

\begin{takeawaybox}
Boundary problems are less about clever algorithms and more about a clean
decomposition with disciplined de-duplication at the seams. Deciding up
front that ``leaves belong to exactly one of the three parts'' is what
keeps the corners from being counted twice.
\end{takeawaybox}
